\section{Interpretability Analysis}
\label{sec:interpretability}

\subsection{Visual Explanation Methods: Grad-CAM and Transformer Attention Rollout}
In clinical dermatology, a black-box malignancy probability without spatial grounding cannot be safely audited at the point of care~\cite{Samek2021explaining,Caruana2015intelligible}. Because our adapted architectures span both convolutional backbones (\texttt{ResNet18Medical}) and hybrid CNN-Transformer networks (\texttt{MedHybridTransformer}), \texttt{src/interpretability.py} implements two complementary attribution mechanisms:

\begin{enumerate}
  \item \textbf{Grad-CAM~\cite{Selvaraju2017grad}:} Given target class score $y^c$ and final spatial feature maps $A^k \in \mathbb{R}^{H'\times W'}$ (from \texttt{layer4[-1]} in \texttt{ResNet-18} or \texttt{stem.stage4[0]} in \texttt{MedHybridTransformer}), channel importance weights are computed via global average pooling of gradients, $\alpha_k^c = \frac{1}{H'W'}\sum_{i,j}\frac{\partial y^c}{\partial A_{ij}^k}$, yielding the non-negative saliency map $L_{\text{Grad-CAM}}^c = \mathrm{ReLU}\!\left(\sum_k \alpha_k^c A^k\right)$ upsampled to $224\times 224$. Forward and backward hooks are cleaned up deterministically inside \texttt{finally} blocks.
  \item \textbf{Transformer Attention Rollout~\cite{Abnar2020quantifying}:} For \texttt{MedHybridTransformer}, we extract the multi-head self-attention matrices $W^{(\ell)} \in \mathbb{R}^{H\times N\times N}$ across all four Workshop~2-adapted Pre-LN layers ($\ell=1\dots 4$), average across heads, account for residual skip connections via $\tilde{A}^{(\ell)} = 0.5\,\bar{W}^{(\ell)} + 0.5\,I$ (row-normalized), and recursively multiply $\hat{A} = \tilde{A}^{(4)}\tilde{A}^{(3)}\tilde{A}^{(2)}\tilde{A}^{(1)}$. Reading the \texttt{[CLS]} row over spatial patch tokens ($14\times 14$) produces a global attention attribution map showing whether long-range Transformer self-attention attends to lesion borders and pigment networks.
\end{enumerate}

\subsection{Quantitative Anatomical Alignment Against Ground-Truth Lesion Masks}
Visual inspection of heatmaps alone is vulnerable to confirmation bias~\cite{Ghassemi2021false}. Using the official Tschandl pixel-level lesion segmentation masks distributed with HAM10000 (\texttt{HAM10000\_segmentations\_lesion\_tschandl.zip}), we validate saliency alignment quantitatively across 15 stratified test cases (5 confident True Positives, 5 highest-risk False Negatives, and 5 highest-confidence False Positives) using four objective metrics:
\begin{itemize}
  \item \textbf{Energy-in-Lesion Ratio:} Fraction of total saliency mass falling inside the ground-truth lesion mask $M$: $\sum_{x\in M} S(x) / \sum_x S(x)$. Across the 15 evaluated test cases, our model achieves a mean Energy-in-Lesion Ratio of \WSixXaiLesionEnergy{}.
  \item \textbf{Pointing Game Accuracy:} Fraction of cases whose peak saliency pixel $\arg\max_x S(x)$ lies strictly inside $M$ (observed: \WSixXaiPointingAcc{}).
  \item \textbf{Saliency-Mask IoU@0.5:} Intersection-over-Union between thresholded saliency $\{S \ge 0.5\}$ and lesion mask $M$ (observed: \WSixXaiIoU{}).
  \item \textbf{Border Shortcut Ratio:} Saliency mass concentrated in the outer 12\% image frame (observed: \WSixXaiBorderEnergy{}), specifically designed to flag shortcut learning on circular dermoscope rims, dark corners, or ruler tick marks.
\end{itemize}

\subsection{Qualitative Evaluation and Five Concrete Failure Modes}
For every selected test case, \texttt{src/interpretability.py} generates a four-panel clinical audit card under \texttt{reports/figures/interpretability/} displaying (1)~the dermoscopic image overlaid with the ground-truth lesion contour and patient metadata (\texttt{age\_group}, \texttt{sex}, \texttt{dx\_type}), (2)~top-3 class probabilities and calibrated malignancy risk relative to the operating threshold, (3)~Grad-CAM localization, and (4)~Workshop~2 Hybrid Attention Rollout. True positive melanomas and basal cell carcinomas consistently localize to asymmetric pigmented networks, atypical globules, and ulcerated vascular cores.

Equally critical for clinical safety, we systematically isolate \textbf{five concrete failure modes} from the test set (Figure~\ref{fig:failure_modes} and Table~\ref{tab:failure_modes}) covering missed melanomas, rare non-melanoma malignancies, pigmented benign keratosis false positives, border/artifact shortcut activations, and subgroup underperformance. Each mode is paired with its patient-harm severity and a concrete clinical verification step allowing a dermatologist or primary-care physician to detect and override the algorithmic error.

\begin{figure}[htbp]
\centering
\safeincludegraphics[0.92\textwidth]{../figures/failure_modes_grid.png}{Visual catalog of 5 concrete failure modes from the HAM10000 test split with Grad-CAM overlays and lesion/border energy diagnostics.}
\caption{\textbf{Five concrete test-set failure modes with Grad-CAM attribution diagnostics.} Inspecting false negatives and false positives exposes morphological mimics and residual border shortcuts where clinicians must override the automated score.}
\label{fig:failure_modes}
\end{figure}

\begin{table}[htbp]
\centering
\caption{\textbf{Structured clinical catalog of five concrete test-set failure modes (\texttt{reports/failure\_modes\_catalog.json}).}}
\label{tab:failure_modes}
\scriptsize
\begin{tabular}{p{1.6cm}p{2.1cm}p{2.3cm}p{3.3cm}p{5.3cm}}
\toprule
\textbf{Mode} & \textbf{Case ID} & \textbf{Error (P(mal))} & \textbf{Root Archetype} & \textbf{Clinician Catch Mechanism} \\
\midrule
\multicolumn{5}{c}{\textit{Pending ROCm interpretability run.}} \\
\bottomrule
\end{tabular}

\end{table}
